\scanpage{74}
\mspage{I}
% archive foliation: 73
Dans le cas où $k$ est de car. $0$, on peut supposer que $k$ est
\deleted{de type fini sur $\mathbf Q$}
le corps des fractions d'un schéma intègre $S$ de type fini sur $\mathbf Z$,
et que $f:X\mathbin{\unsure{-|\to}}Y$ provient d'un morphisme de
$S$-préschémas $\overline f:\overline X\mathbin{\unsure{-|\to}}\overline Y$,
intègres, $\overline Y$ simple sur $S$, et $\overline X$
\deleted{simple} normal sur $S$, les fibres de l'un et l'autre étant
\unsure{géomt} connexes.
\deleted{Alors \unsure{pour}
$\forall\,\unsure{y\in Y|s\in S}$ \unsure{sc}
$H^1(Y,\underline O_Y)\to H^1(X,\underline O_X)$
est de dim $r$, il existe un
\unsure{peut supposer qu'il est constant}.}
On peut supposer aussi $X$ et $Y$ coh. plats sur $S$ en dim $1$,
\deleted{\unsure{enfin que}}
Alors \deleted{\unsure{celui}}
l'hom
$R^1q_*(\underline O_{\overline Y})\mathbin{\unsure{-|\to}}
R^1p_*(\underline O_{\overline X})$
est injectif, \deleted{et que}
(car il l'est en tt pt par \unsure{évidence|réduction}).
Alors \unsure{rétrécissant} $S$, si le \unsure{conoyau}
du $H^1(Y,\underline O_Y)\mathbin{-}H^1(X,\underline O_X)$
est de dim $r$, la $=$ chose sera vraie pour les fibres en \uline{tous}
les points $s\in S$, en particulier en des pts de
